\cotefile{137}
\scanpage{31}
(Si $\mathcal A$ est réduit à un objet, \inserted{$A$ \unsure{soit|seul} $k$} l'anneau
des endom $k$ de celui-ci, on
trouve une équivalence avec la
cat. des $k$-modules à droite
libres de type fini, par $X\to\operatorname{Hom}(A,X)$)

Transformant par dualité les énoncés.
Si tt objet de $C$ est $\varprojlim$
d'objets de $\mathcal A$, alors le foncteur
canonique
\[
 v^{\circ}\colon C^{\circ}\longrightarrow \widehat{\mathcal A^{\circ}}_{\unsure{\mathrm{Ens}|\mathrm{Et}}}
 \quad(\text{resp. }\widehat{\mathcal A^{\circ}}_{\mathrm{ab}})
\]
vers les foncteurs $\underline{\operatorname{Hom}}(\mathcal A,\mathrm{Ens})$ (resp. $\underline{\operatorname{Hom}}(\mathcal A,\mathrm{Ab})$)
\uline{covariants} sur $\mathcal A$, à valeurs dans
$(\mathrm{Ens})$ (resp. $\mathrm{Ab}$) est pl. fidèle. Si
de plus, tt objet de $C$ est
$\unsure{\underline{\lim}|\varprojlim}$ de certains \unsure{types|type} d'objets
de $C$, alors l'image essentielle est
contenue dans la sous-cat. pleine
de $\widehat{\mathcal A^{\circ}}_{\mathrm{Ens}}$ (resp. $\widehat{\mathcal A^{\circ}}_{\mathrm{ab}}$) formée des
foncteurs $\mathcal A\to\mathrm{Ens}$ (resp. $\mathcal A\to\mathrm{Ab}$,)
qui sont $\varprojlim$ des dits types des
foncteurs représentables -- et
l'image essentielle est $=$ égale à
cette dernière catégorie si
% archive sheet number: 30
\scanpage{32}
$C$ est stable sous les limites projectives
de type envisagé.

En particulier, on trouve

\uline{Prop} Soient $C$ une cat. (resp. cat
additive) $\mathcal A$ une sous-cat. pleine,
$\mathcal T,\mathcal T'$ deux sous-\deleted{catégories pleines}\inserted{ens}
de $\operatorname{Ob}(\mathrm{Cat})$, \unsure{suppos.|supposons} que $C$ soit
stable par $\varprojlim$ de type $\mathcal T$ et $\varinjlim$
de type $\mathcal T'$, et que tt objet de
$C$ soit une $\varprojlim$ \inserted{dans $C$} de type $\mathcal T$ d'objets
de $\mathcal A$, et une $\varinjlim$ dans $C$ de type $\mathcal T'$
d'objets de $\mathcal A$. Alors $C$ est équivalente
à la fois à la sous cat. pleine
de $\widehat{\mathcal A}$ \inserted{$=\underline{\operatorname{Hom}}(\mathcal A^{\circ},\mathrm{Ens})$}
(resp. $\widehat{\mathcal A}_{\mathrm{ab}}$ \inserted{$=\underline{\operatorname{Hom}}(\mathcal A^\circ,\mathrm{Ab})$})
formée des
foncteurs \inserted{$\mathcal A^\circ\to\mathrm{Ens}$ (ou $\mathcal A^\circ\to\mathrm{Ab}$)} qui sont $\varprojlim$ \deleted{de type $\mathcal T$}
de type $\mathcal T$ de foncteurs repr,
et \unsure{à|de} la \inserted{cat. opposée de la} sous-catégorie pleine de
$\widehat{\mathcal A^\circ}\simeq\underline{\operatorname{Hom}}(\mathcal A,\mathrm{Ens})$ (resp. $\widehat{\mathcal A^\circ}_{\mathrm{ab}}\simeq
\underline{\operatorname{Hom}}(\mathcal A,(\mathrm{Ab}))$) qui sont \deleted{stables par}\inserted{\unsure{respect$^{\mathrm{t}}$|resp.} des}
$\unsure{\underset{\longleftrightarrow}{\lim}|\varprojlim|\varinjlim}$ de \deleted{type} $\mathcal T'$ de foncteurs
représentables.
% archive sheet number: 31
\scanpage{33}
En fait, on trouve un \unsure{diagramme|diagr.me|diagr.}
ess. commutatif
\diagram{\begin{tikzcd}[ampersand replacement=\&,column sep=0.45em,row sep=3em]
\&[-1em]\mathcal A \arrow[dl,"{\underline{\mathrm{Id}}}"'] \arrow[drrr,no head,"{\underline{\mathrm{id}}}"'] \& \subset \& C \arrow[dll,"{\simeq}"',"{v}"] \arrow[drr,no head,"{v'}","{\simeq}"'] \& \& \& \\
\mathcal A \& {\subset\widehat{\mathcal A}_{\mathcal T}} \& {\subset\widehat{\mathcal A}} \& \& \mathcal A \arrow[r,hook] \& {(\widehat{\mathcal A^\circ}_{\mathcal T'})^\circ} \arrow[r,hook] \& {(\widehat{\mathcal A^\circ})^\circ}
\end{tikzcd}}
(resp. en $\widehat{\mathcal A}_{\mathrm{ab}}$ et $\widehat{\mathcal A^\circ}_{\mathrm{ab}}$)

On trouve en particulier une
équivalence remarquable
\[
 \widehat{\mathcal A}_{\mathcal T}\simeq
 (\widehat{\mathcal A^\circ}_{\mathcal T'})^\circ
\]
qui peut se définir intrinsèquement \inserted{en termes de $\mathcal A,\mathcal T,\mathcal T'$,} chaque fois qu'on a une
sous-cat. $D$ de $\widehat{\mathcal A}$ (resp. une sous cat.
additive de $\widehat{\mathcal A}_{\mathrm{ab}}$), stable \inserted{dans $\widehat{\mathcal A}$} par
$\varprojlim$ de type $\mathcal T$, et \inserted{où existent les} $\varinjlim$ de type $\mathcal T'$,
(pas néc. les $=\mathrm{s}$ que dans $\widehat{\mathcal A}$ !)
et dont chaque objet est à la fois
$\varprojlim$ \inserted{de type $\mathcal T$} dans $C$ de foncteurs repr.,
et $\varinjlim$ \inserted{de type $\mathcal T'$} dans $C$ de tels foncteurs.

Si $\mathcal A$ est réduit à un seul objet, \inserted{$A$} et
dans le cas \unsure{resp$^\circ$|resp.} ($C$ additive)
$\mathcal T=\mathcal T'=$ ens. catégories discrètes finies,
on trouve une antiéquivalence
entre modules libres de t. f à g. et
% archive sheet number: 32
\scanpage{34}
à droite sur $k=\operatorname{End}(A)$, et ladite
antiéquivalence n'est autre que la
dualité de tels modules.

Revenons \deleted{\unsure{notre}}\inserted{à la} condition où
tt objet de $C$ est une $\varinjlim$ \unsure{dans|de} $C$
de type $\mathcal T'$ d'objets de $\mathcal A$. Alors
\unsure{tte|pr|toute|si} catégorie $C'$ où ces limites existent,
si $\underline{\operatorname{Hom}}_{\mathcal T'}(C,C')$ est la sous-catégorie
pleine de $\underline{\operatorname{Hom}}(C',C')$ formée des
foncteurs qui commutent à ces
$\varinjlim$, on trouve que
\[
 \underline{\operatorname{Hom}}_{\mathcal T'}(C,C')\longrightarrow\underline{\operatorname{Hom}}(\mathcal A,C')
\]
% slanted
\margin{\uline{à vérifier}}
est pl. fidèle \deleted{(et même une équival.
de cat. si \illegible{2} $C$ est libre}

Supposons $C$ additive, et $\mathcal T'$ l'ens
des cat. discrètes finies -- donc
tt objet de $C$ est somme d'objets
de $\mathcal A$, alors on trouve, si $C'$ additive
\[
 \underline{\operatorname{Hom}}_{\mathrm{add}}(C,C')\simeq
 \underline{\operatorname{Hom}}_{1/2\,\mathrm{add}}(\mathcal A,C')
\]
\unsure{Se rappelant|Si, rappelons|Si on applique} de l'inclusion
\[
 C\lhook\joinrel\xrightarrow{\mathrm{pl.\ fid.}}
 \deleted{\widehat{\mathcal A}_{\mathrm{ab}}}
 \underline{\operatorname{Hom}}_{1/2\,\mathrm{add}}(\mathcal A^\circ,\mathrm{Ab})
\]
et désignant par $D\subset\deleted{\widehat{\mathcal A}_{\mathrm{ab}}}\unsure{''}$ l'image
essentielle, on trouve alors
% archive sheet number: 33
\scanpage{35}
\[
 \underline{\operatorname{Hom}}_{\mathrm{add}}(D,C')\xrightarrow{\simeq}
 \underline{\operatorname{Hom}}_{1/2\,\mathrm{add}}(\mathcal A,C')
\]
\unsure{d'où|allant} en sens inverse
\[
 \underline{\operatorname{Hom}}_{1/2\,\mathrm{add}}(\mathcal A,C')\xrightarrow{\simeq}
 \underline{\operatorname{Hom}}_{\mathrm{add}}(D,C')
\]
d'où un accouplement
\diagram{\begin{tikzcd}[ampersand replacement=\&,row sep=0pt,column sep=0.6em]
D \& \times \& \underline{\operatorname{Hom}}_{1/2\,\mathrm{add}}(\mathcal A,C') \arrow[r] \& C' \\
\rotatebox{-90}{$\subset$} \& \& \& \\
\underline{\operatorname{Hom}}_{1/2\,\mathrm{add}}(\mathcal A^\circ,\mathrm{Ab}) \& \& \&
\end{tikzcd}}
qu'on note par \deleted{\unsure{la|ex.}}
\[
 (F,G)\unsure{\mapsto|\vdash} F\underset{C}{\otimes}G
\]
[dans le cas où $\mathcal A$ est réduit à un
élément, donc où sa str. $1/2$ add. est
\unsure{défini|définie} par un anneau $k$,
$F$ est un \inserted{$k$-}module à droite libre de
t. f, $G$ est un \deleted{\unsure{$k$}}objet de $C'$ où $k$
opère à g (i.e. un $k$-$C'$-module
à gauche) et $F\otimes G$ est le
produit tensoriel ``ordinaire''.

Si on considère \deleted{\unsure{ce|alors}}le plongement
dual $C\xrightarrow{\sim}D'{}^\circ\subset\underline{\operatorname{Hom}}_{1/2\,\mathrm{add}}(\mathcal A,\mathrm{Ab})^\circ$,
on trouve
\diagram{\begin{tikzcd}[ampersand replacement=\&]
\underline{\operatorname{Hom}}_{\mathrm{add}}(D'{}^\circ,C') \arrow[r,"{\simeq}"] \&
\underline{\operatorname{Hom}}_{1/2\,\mathrm{add}}(\mathcal A,C') \arrow[l,bend left=40,"{\simeq}"]
\end{tikzcd}}
d'où
% archive sheet number: 34
\scanpage{36}
i.e.
\diagram{\begin{tikzcd}[ampersand replacement=\&,row sep=1.3em,column sep=0.8em]
D'{}^\circ \arrow[d,hook] \& \times \& \underline{\operatorname{Hom}}_{1/2\,\mathrm{add}}(\mathcal A,C') \arrow[r] \& C'\\
\underline{\operatorname{Hom}}_{1/2\,\mathrm{add}}(\mathcal A,\mathrm{Ab})^\circ \& \& \&
\end{tikzcd}}
noté
\[
 \underline{\operatorname{Hom}}_C(F,G)
\]

-\, -\, -\, -

\noindent\hfil\rule{0.4\linewidth}{0.35pt}\hfil

On suppose maintenant qu'on a
une catégorie $\mathscr S$, une sous-catégorie
\inserted{pl $C$,} de $\mathscr S_{\mathrm{ab}}$ stable par sommes directes
(alors $C$ est additive) et une
sous-catégorie pleine $\mathcal A$ de $C$ telle
que tt objet de $C$ soit isom à une
somme directe \inserted{dans $C$ donc dans $\mathscr S_{\mathrm{ab}}$} \deleted{\unsure{de}} d'objets de $\mathcal A$.

On a alors \deleted{\unsure{fonct.}}
\[
 \unsure{\alpha|u},\quad \mathscr S^\circ\longrightarrow
 \deleted{\widehat{\mathscr S}_{\mathrm{ab}}}
 \underline{\operatorname{Hom}}_{\mathrm{add}}(C,\mathrm{Ab})
\]
\deleted{\unsure{comme}} en associant à tt $X\in\operatorname{Ob}\mathscr S$
le foncteur $M\longmapsto\operatorname{Hom}_{\mathscr S}(X,M)$.
Or ce foncteur est connu à équiv.
près par \unsure{de --|sa --|quand on} \unsure{connaissant|en connaissant|connaît} son image
dans $\underline{\operatorname{Hom}}_{\unsure{.\mathrm{add}|\mathrm{add}|1/2\,\mathrm{add}}}(\mathcal A,(\mathrm{Ab}))$, donc l'étude
de $\alpha$ est \unsure{équivalente} à celle de
% archive sheet number: 35
\scanpage{37}
\[
 \mathscr S^\circ\longrightarrow\underline{\operatorname{Hom}}_{1/2\,\mathrm{add}}(\mathcal A,\mathrm{Ab})
\]
Nous ferons cette étude dans le
cas où $\mathscr S=\widehat\Delta=\mathrm{Ss}$, et dans
chacun des 2 cas suivants
\deleted{\unsure{en supposant qu'on a}}

$1^\circ$) $C$ est formé des groupes ab. $1/2$ simpliciaux
\deleted{\unsure{modulaires|multiplicatifs}}\inserted{\unsure{str.|ab.} $\partial$-parfaits}, $\mathcal A$ est la
% slanted
\margin{$C$ abélienne}
sous catégorie formée des
$\Phi(i,n)$ ($n\in\mathbb N^*$, \deleted{$\unsure{j\geq1}$} $i\geq1$) et $\Lambda^j\Psi_*$ ($j\geq0$).

On note que \inserted{pour $X\in\operatorname{Ob}\widehat\Delta$ fini,} la connaissance de
$\operatorname{Hom}_{\widehat\Delta}(X_{\unsure{\bullet|*}},M_*)$ comme foncteur
covariant en $M_*$ est réduit à
celle de $\operatorname{Hom}_{\widehat\Delta}(x_{\bullet},\deleted{\unsure{F}}\varphi)$ pour
$\varphi$ de $\mathcal A$. On peut encore
réduire \deleted{\unsure{aux}}\inserted{pour} $\varphi$ de $\mathcal A_0$ défini
ci-dessous --, et $=$ pour $\varphi$ parcourant
les $\Phi^i$.

$2^\circ$) $C_0$ \inserted{$\subset C$} \deleted{est} formé \unsure{:} des groupes ab
% slanted
\margin{$C_0$ est additive mais non abélienne.}
$1/2$ simpliciaux \deleted{\unsure{des}} spéciaux
(\deleted{\unsure{dans}}\inserted{i.e.} $\in C$ et les $\pi_i(M_{\bullet})$ sont
sans torsion), $\deleted{\unsure{\Phi}}\mathcal A_0\subset C_0$ sous cat.
de ceux qui sont de la forme
$\Phi^{\unsure{0|i}}=\Lambda^i\Phi_*\deleted{\unsure{\Psi}}$ ($i\geq1$) ou $\Psi^j=\Lambda^{\unsure{0|j}}\Psi_*$ ($j\geq0$).
% archive sheet number: 36
\scanpage{38}
Il y a lieu de \unsure{s'intéresser|travailler}
$1^\circ$) au cas où on a un anneau
de base $k$ de Dedekind,
et $2^\circ$) \unsure{au|un|avec} $k$ quelconque \inserted{(\unsure{évt comm})}
(donc $\unsure{\Theta|\mathcal C|C}\subset\widehat\Delta_{k\text{-}\mathrm{mod}}$). La \deleted{\unsure{distinction|situation}}
conclusion à retenir, c'est
qu'on connaît les Hom de
$X_*$ dans tt $k$-module str.
$\partial$-\deleted{\unsure{parfaits}}\inserted{parfait $M$ (comme foncteur de $M$)} dès qu'on connaît
le complexe \inserted{de cochaînes ``non dég''} $C^{\unsure{*!|*}}(X_*,k)$ à coeff. dans
$k$. Notre propos principal
est de voir comment reconstituer
la structure multiplicative
et ``puiss. divisées'') sur le
foncteur
\[
 M\mapsto F_{X_*}(M_*)=F(M_*)=\operatorname{Hom}_{\widehat\Delta}(X_*,M_*)
\]
i.e. les homs
\[
 \left\{\begin{aligned}
 F(M_*)\otimes_k F(N_*)&\longrightarrow F(M_*\otimes N_*)\\
 \Gamma^i F(\deleted{\unsure{\Gamma^i}}M_*)&\longrightarrow F(\Gamma^i(M^{\unsure{n|*|\bullet}}))
 \end{aligned}\right.
\]
en tant que \deleted{\illegible{1}} structures
sur le complexe \inserted{de cochaînes} $C^*=C^{*!}(X_*,k)$

Il faudrait examiner si pour
$k$ de car 0 \inserted{$C^*$ avec} cette structure suppl.
% archive sheet number: 37
\scanpage{39}
est déterminé par le complexe de
De Rham-Sullivan de $X_*$, et dans
le cas $k$ quelconque (et cas $2^\circ$ disons
-- $k$ principal est \unsure{le $1^\circ$|cas $1^\circ$|que $1^\circ$}) il est
déterminé par le complexe de
De Rham à p. d. \deleted{Enfin,}
\deleted{\unsure{il y aura|il y a}} (il est clair que $C^*$
avec les structures suppl. envisagées
\deleted{définit}\inserted{redétermine} le complexe de De Rham-Sullivan resp. le complexe de
De Rham à p. d.) Enfin, il
faut voir dans quelle mesure $C^*$
avec ses structures multiplicatives
à p. d. permet de reconstituer le
type d'homotopie du $X_*$ dont
il provient. (Enfin, bien sûr,
\unsure{faisceautiser|faisceauter} et relativiser -- !!!)
% archive sheet number: 38
\scanpage{40}
Soient \inserted{$k$ anneau comm.} $\unsure{\mathcal C|\mathcal E|\mathbb K}$ une catégorie \inserted{$k$-}additive (i.e. additive, munie
de $k\to\operatorname{End}(\mathrm{id}_{\mathcal C})$), $C^\bullet$ un complexe \inserted{de cochaînes} dans
$\mathcal C$
\[
 \cdots\longrightarrow C^0\xrightarrow{d^0}C^1\xrightarrow{d^1}
 \cdots\longrightarrow C^i\xrightarrow{d^i}C^{i+1}\longrightarrow\cdots
\]
On suppose :

a)
\[
\operatorname{Hom}(C^i,C^j)\simeq
\begin{cases}
0&\text{si }j\ne i,i+1\\
k&\text{si }j=i,\ \text{base }\mathrm{id}_{C^i}\\
k&\text{si }j=i+1,\ \text{base }d^i
\end{cases}
\]

b) Tout objet de $\mathcal C$ est isomorphe à une
somme \inserted{finie} d'objets de la forme $C^i$.

Soit $\mathcal A$ la sous-catégorie pleine de $\mathcal C$
formée des \deleted{$\unsure{\Phi}$} $C^i$ ; elle n'est pas additive,
mais hérite d'une structure $k$-linéaire,
i.e les $\operatorname{Hom}_{\mathcal A}$ sont des $k$-modules, leur
composition est $k$-bil.

\deleted{\unsure{C|E}} L'existence de $(\mathcal C,C^\bullet)$ satisfaisant a) et
b) est immédiate par construction ``abstraite''.
On peut aussi réaliser $\mathcal C$ comme \inserted{la} cat. \inserted{$K_{\bullet}$} des
complexes de chaînes \inserted{$L_{\bullet}$} de $k$-modules
satisfaisant les conditions

$1^\circ$) \deleted{\unsure{$K_{\bullet}$}} $L_{\bullet}$ est acyclique i.e. $H_i(L_{\bullet})=0$ $\forall i\geq0$

$2^\circ$) $L_n$ et $Z_n(L_{\bullet})$ sont libres de t. f. pour tt $n$

$3^\circ$) $L_n$ est nul pour $n$ grand.

On prendra alors :
\[
 C^i=\deleted{\unsure{K^i_{\bullet}}} C^i_{\bullet}=
 \{\,0\longrightarrow\cdots\longrightarrow0\longrightarrow
 \overset{\text{deg }i+1}{k}\xrightarrow{\mathrm{id}}
 \overset{\text{deg }i}{k}\longrightarrow0\longrightarrow\cdots\longrightarrow
 \overset{\text{deg }0}{0}\,\}
\]
i.e. les comp. sont nulles sauf en degré $i,i+1$ où on
aura $k$, et l'app. diff. \deleted{\unsure{non nulle}}\inserted{$i+1$} est l'identité.
On prend $d^i\colon \deleted{\unsure{K}}C^i_{\bullet}\xrightarrow{d^i_{\bullet}}
\deleted{\unsure{K}}C^{i+1}_{\bullet}$ caractérisé
par $(d^i)_{i+1}=\mathrm{id}_k$. Les conditions a) et b) sont
satisfaites trivialement. On pourrait aussi
prendre des complexes de cochaînes à la place,
% margin note written sideways
\margin{avec les $C_i^\bullet=(\overset{\text{deg }0}{0}\to\cdots\to0\to
\overset{\text{deg }i}{k}\to\overset{\text{deg }i+1}{k}\to\cdots$, mais
ici on a $C_i^\bullet\longleftarrow C_{i+1}^\bullet$, \deleted{\unsure{ce}} il faut en prendre
la cat. \uline{opposée} de celle \inserted{des} cochaînes \unsure{en quest.|envisagée}.
Mais on le retrouvera en \unsure{2 repr.|? repr.|2 vers.} de façon plus
\uline{élégante} dans les Ex.}
% archive sheet number: 39
\scanpage{41}
\deleted{\{ $\unsure{\mathcal C|\mathcal E}$. catégorie \{ des $k$-modules $1/2$ simpliciaux \deleted{\unsure{tels}}\inserted{$L_*$} qui
sont flasques et str. $\partial$-parfaits, i.e. qui satisfont
deux \deleted{aux} conditions équivalentes

a) $L_*$ est \deleted{\unsure{équivalent}}\inserted{isomorphe} à une somme directe finie
de modules $1/2$ simpliciaux
$\deleted{\unsure{\Phi^i_*=\Lambda^i\Phi_*|\Phi(i)_*=\Lambda^i\Phi_*}\ (i\geq1)}\inserted{C^i_*=\Lambda^{i+1}\Phi_{*k}\ (i\geq0)}$
}
% slanted
\margin{\deleted{On suppose tt module proj. de type fini sur $k$ est libre p. ex. $k$ principal}}
\deleted{b) Les $L_n$ sont libres de type fini, de rang constant
sur $\operatorname{Spec} k$, la fonction $n\mapsto\operatorname{rg}_k L_n$ est polynomiale
pour $n$ grand, $L_*$ est flasque i.e. $\pi_i(L_*)=0$ pour $i\geq0$.

$\mathcal A\subset\mathcal C$ la sous-catégorie pleine formée des
$\deleted{\unsure{\Phi^i_*}}=C^{i-1}_*$

\uline{NB} $\mathcal C$ est additive (et karoubienne), $\mathcal A$ n'est pas
stable par sommes finies, mais hérite une structure
$k$-linéaire. On a
\[
 \operatorname{Hom}(\deleted{\unsure{\Phi}}C^i_*,\deleted{\unsure{\Phi}}C^j_*)=
 \begin{cases}
 0&\text{si }j\ne i,i+1\\
 k&\text{si }j=i,\ \text{base }\mathrm{id}_{\unsure{\Phi(i)|C^i|C^i_*}}\\
 k&\text{si }j=i+1,\ \text{base }d_i=T_*\wedge
 \end{cases}
\]
Notons que
\deleted{$\unsure{\Phi_*}\ \unsure{\Phi^i_*|C^i_*}=C^{i+1}_*$ \unsure{inversement} $C^i_*=\unsure{\Phi^{i+1}_*}$ $(i\geq0)$}

$C^\bullet_*=(C^i_*,T_*\wedge)$ est un complexe de cochaînes dans
$\mathcal A$ ou dans $\mathcal C$ ou dans $\widehat\Delta_{k\text{-}\mathrm{lin}}$ -- il jouera
le rôle du complexe de cochaînes \uline{universel}.}

\deleted{\unsure{En effet}}

\uline{Prop} 1. Soit $\mathcal B$ une catégorie \deleted{additive}\inserted{\deleted{\unsure{lin.}} $k$-additive}. Alors
\diagram{\begin{tikzcd}[ampersand replacement=\&,column sep=2em,row sep=1.4em]
\underline{\operatorname{Hom}}_{\deleted{\unsure{\mathrm{add}|k\text{-}\mathrm{lin}}}\unsure{k\text{-}\mathrm{lin}}}(\unsure{\mathcal K|\mathcal C},\mathcal B)
 \arrow[r,"{\substack{\mathrm{rest}\\\simeq}}"] \&
 \underline{\operatorname{Hom}}_{k\text{-}\mathrm{lin}}(\mathcal A,\mathcal B)
 \arrow[d,"{\simeq}"',"{F\mapsto F(C^\bullet_*)}"] \\
\& {\mathcal B^\bullet_{\unsure{k}}=\text{cat. des \inserted{complexes de} cochaînes de }\mathcal B}
\end{tikzcd}}

\deleted{\uline{\unsure{Question|Corollaire}} Soit $K^\bullet\in\operatorname{Ob}\mathcal B^*$ un complexe de cochaînes de $\mathcal B$}

\uline{Exemple 1} \inserted{faisons $\mathcal B=\mathrm{Ab}_k$.} Soit $L_*\in\operatorname{Ob}\mathcal C$, considérons le
foncteur cov. qu'il représente
\[
 \mathcal K\longrightarrow\mathrm{Ab}_k:\qquad M_*\longmapsto\operatorname{Hom}_{\mathcal C}(L_*,M_*)
\]
on a ainsi un foncteur pl. fid. (sous cat. des fonct. repr. --)
\diagram{\begin{tikzcd}[ampersand replacement=\&,row sep=0.9em,column sep=1em]
\mathcal K^\circ \arrow[r,hook] \& \underline{\operatorname{Hom}}_{k\text{-}\mathrm{lin}}(\mathcal K,\mathrm{Ab}_k)
 \arrow[dr,"{\simeq}"] \& \& \\
\& \& \underline{\operatorname{Hom}}_{k\text{-}\mathrm{lin}}(\mathcal A,\mathrm{Ab}_k) \arrow[r,"{\simeq}"] \& \mathrm{Ab}^{\bullet\deleted{\unsure{*}}}_{k\deleted{\unsure{*}}}
\end{tikzcd}}
% archive sheet number: 40
\scanpage{42}
On trouve donc un foncteur pl. fid
\[
 \mathcal C^\circ\lhook\joinrel\xrightarrow{\deleted{\unsure{\beta}}\alpha^\bullet}
 \mathrm{Ab}_k^{\bullet\deleted{\unsure{*}}}
\]
qui associe à tout $\deleted{\unsure{M_*}}L_*$ associe le complexe
de cochaînes $\operatorname{Hom}_{\mathcal C}(L_*,C^\bullet_*)$
\[
 0\longrightarrow\operatorname{Hom}(L_*,C^0_*)\longrightarrow
 \operatorname{Hom}(L_*,C^1_*)\longrightarrow\cdots,
\]
et dont la restriction à $\mathcal A$ est donnée par
\[
 \alpha^\bullet\deleted{\unsure{\beta}}(C^i_*)\longmapsto
 (0\to0\to\cdots\,0\to\underset{\text{degré }i}{k}
 \xrightarrow{\mathrm{id}}\underset{\text{degré }i+1}{k}\to0\to\cdots)
 =\deleted{\unsure{K(i)}}\inserted{C_i^\bullet}
\]
et par la loi
\diagram{\begin{tikzcd}[ampersand replacement=\&,column sep=1.2em,row sep=0.65em]
\& \alpha^\bullet(C^i_*) \& \alpha(C^{i+1}_*) \arrow[l,"{\deleted{\unsure{\beta}}\alpha^\bullet(d^i)}"'] \&\\
\& \Vert \& \Vert \&\\
\& \deleted{\unsure{K}}C_i^\bullet \arrow[r] \& \deleted{\unsure{K}}C_{i+1}^\bullet \& k-0\\
\& \vdots \& \vdots \&\\
\& 0 \arrow[d] \& 0 \arrow[d] \&\\
i \& k \arrow[d] \& 0 \arrow[l] \arrow[d] \&\\
i+1 \& k \arrow[d] \& k \arrow[l,"{\mathrm{id}}"'] \arrow[d] \&\\
i+2 \& 0 \arrow[d] \& k \arrow[l] \arrow[d] \&\\
\& 0 \& 0 \&\\
\& \vdots \& \vdots \&
\end{tikzcd}}
\margin{L'image essentielle \inserted{\unsure{pl.|$K^i$}} de $\mathcal C^\circ$ est la catégorie des $K^\bullet\in\operatorname{Ob}\mathrm{Ab}_k^\bullet$ qui sont libres de t. f en tte dimension, et qui sont à coh. nulle (en admettant que tt facteur direct d'un libre de t. f. est itou}
% slanted
\margin{\unsure{Commencer|Commençons|commentaire} par Ex 2}

\uline{Ex 2} Faisons $\mathcal B=\mathrm{Ab}_k^\circ$,
qui est $k$-linéaire (on suppose
$k$ commutatif), et pour tt $L_*\in\operatorname{Ob}\mathcal C$, on considère le foncteur contrav. qu'il représente, interprété
comme foncteur $k$-lin.
\[
 \mathcal C^\circ\longrightarrow\deleted{\unsure{\mathrm{Ab}^\circ_k}}\mathrm{Ab}_k^{*}
 \qquad M_*\mapsto\operatorname{Hom}_{\mathcal C}(M_*,L_*)
\]
on trouve ainsi un foncteur pl. fid (par la
\unsure{même voie|méthode|même idée|covariant}
\diagram{\begin{tikzcd}[ampersand replacement=\&,column sep=1em,row sep=1.5em]
\mathcal C \arrow[r,hook] \& \underline{\operatorname{Hom}}_{k\text{-}\mathrm{lin}}(\mathcal C^\circ,\mathrm{Ab}_k^\bullet) \arrow[r,"{\sim}"] \arrow[d,"{\simeq}"] \& \underline{\operatorname{Hom}}_{k\text{-}\mathrm{lin}}(\mathcal C,\mathrm{Ab}_k^\circ)^\circ \arrow[d,"{\simeq}"]\\
\& \underline{\operatorname{Hom}}_{k\text{-}\mathrm{lin}}(\mathcal A^\circ,\mathrm{Ab}_k^\bullet) \arrow[r,"{\sim}"] \arrow[d] \& \underline{\operatorname{Hom}}_{k\text{-}\mathrm{lin}}(\mathcal A,\mathrm{Ab}_k^\circ)^\circ \arrow[d,"{\simeq}"]\\
\& {(\mathrm{Ab}_k^\circ)_{\deleted{\unsure{*}}}^{*}\deleted{\simeq(\mathrm{Ab}_{k*}^{\unsure{*}})^\circ}} \& {(\mathrm{Ab}_k^{\circ *})^\circ}
\end{tikzcd}}
ou encore un foncteur pl. fid.
\[
 \mathcal C^{\unsure{\bullet|\circ}}\lhook\joinrel\xrightarrow{\beta_{\bullet}}\mathrm{Ab}_{k*}
\]
% archive sheet number: 41
\scanpage{43}
On trouve ici que
\[
 \beta_\bullet(C^i)=\deleted{\unsure{C^i_\bullet}}
 (\ldots 0\to\ldots 0\to\overset{\deleted{\unsure{i-1|i|i+1}}i}{k}
 \to\overset{\deleted{\unsure{i|i+1}}i+1}{k}\to0\to\ldots\to0)
 \quad\inserted{\text{pour }i\geq1}.
\]
\inserted{$\beta_\bullet(C^0)=(\ldots\deleted{\underset{\text{deg }0}{k}}\to0\to
 \underset{\text{deg }0}{k}\to0)$}
et $\beta_\bullet(d^i): C^i_\bullet\xrightarrow{d^i_\bullet} C^{i+1}_\bullet$
est donné par
\[
 (d^{\deleted{\unsure{i+1}}i}_{\deleted{\unsure{i+1|i}}\bullet})=\mathrm{id}_k
 \qquad \unsure{i\geq10|i\geq0}
\]
\deleted{$d\qquad\unsure{i-1|i=0}$}

% slanted
\margin{Ce n'est pas bien bon à cause de l'irrégularité pour \unsure{ces couples} plutôt}
\deleted{On \unsure{donne|a donc} \unsure{comme|ainsi} \unsure{type|un} modèle la cat.
(\unsure{des} \inserted{$/k$}) complexes de chaînes \unsure{déj. conn.} que nous
\unsure{avions|considéré} comme possibilité de construction d'un couple $(C,C^\bullet)$.}

\uline{Passage de l'exemple 1 à l'exemple 2}

Considérons le foncteur
\[
\begin{aligned}
 M&\longmapsto M^\vee\\
 \mathrm{Ab}_k^\circ&\unsure{\longrightarrow|-}\mathrm{Ab}_k
\end{aligned}
\]
qui induit une équivalence de la catégorie des \inserted{$k$-}modules libres de t.f. avec elle-m.
Cela donne aussi une équivalence entre les cat de complexes de chaînes et cochaînes
à \unsure{comp.} \deleted{\unsure{$k^\vee$ de t.f.}} libres de t.f., et en
celles où les $Z_n$ (et $Z^i$) sont libres de t.f. et \unsure{h. libres|les bords}.
\deleted{\unsure{sont libres, en sont tous}. On \unsure{peut}, \unsure{ce qui du point de vue naturel}}
\deleted{\unsure{pour (pos)}} \unsure{a donc} \unsure{le produit}
\inserted{$\alpha_{\bullet}(\unsure{K|x|\mathbb K})\overset{\unsure{\text{défn}|\text{de}|\text{dual}}}{=}\alpha^\bullet(x)^\vee$} \deleted{\unsure{donné}}
un diagramme d'équivalences

\diagram{\begin{tikzcd}[ampersand replacement=\&]
 \deleted{\unsure{K_\bullet|\mathbb K_\bullet}}\& C\arrow[r,"{\alpha_\bullet}","\sim"']\arrow[d,"{(\alpha^\bullet)^\circ}"',"\sim"] \& K_\bullet\\
 \& K^{\bullet\circ}\arrow[ur,dash,"\sim"] \&
\end{tikzcd}}

% margin note written sideways
\margin{donc $\alpha_n(x)\simeq\mathrm{Hom}(C^{n-1},x)$ si $n\geq1$
et les hom $\alpha_n(x)\to\alpha_{n-1}(x)$ pour $n\geq2$
sont donnés par transposition à partir de $C^{n-1}\to C^n$.
Mais $\alpha_0(x)\simeq\mathrm{Hom}(x,C^0)^\vee$
et $\alpha_1(x)\to\alpha_0(x)$ donné par
$\mathrm{Hom}^{\unsure{\vee}}(C^0,x)\to\mathrm{Hom}(x,C^0)^\vee$ \unsure{ce qu'on aura|qu'on devine}.}

\deleted{Il suffit de définir un isom pour les restrictions à $A$ ce qui est trivial.}

\uline{Explicitons \unsure{quelque peu}.}

\uline{Prop 2} On a une dualité parfaite fonctorielle en $M_\bullet\in\mathrm{Ob}\,K_\bullet$
(\unsure{à|de} $n\geq0$ fixé)
\[
\begin{aligned}
 \alpha_n(\unsure{M|x})_{\unsure{\bullet}}\ \times\ 
 \alpha^n(M)&\longrightarrow k\\[-2pt]
 \rotatebox{90}{$=$}\qquad \rotatebox{90}{$=$} &\qquad \rotatebox{90}{$\simeq$}\\
 \mathrm{Hom}(C^n,x)\times\mathrm{Hom}(x,C^{n+1})&\longrightarrow
 \mathrm{Hom}(C^n,C^{n+1})\deleted{\unsure{\to}}\simeq k\\
 u,v&\longmapsto\unsure{vu}.
\end{aligned}
\]
% archive sheet number: 42
\scanpage{44}
Revenant à la prop. 1. Soit $C^\bullet$ \inserted{$\in\mathrm{Ob}\,\mathcal B^\bullet$}
un complexe de cochaînes de $\mathcal B$, comment réconstituer le foncteur $k$-lin
\[
 F_{C^\bullet}: \unsure{\mathcal K|\mathcal C|\mathbb K}\longrightarrow\mathcal B
\]
i.e. comment ``calculer'' $F_C(x)$ pour $x\in\mathrm{Ob}\,K$ ?
Il faut qu'on dise comment on ``repère'' $x$,
disons par $\alpha^\bullet(x)\in\mathrm{Ab}_k^\bullet$. Ceci dit, on trouve

\uline{Prop. 3} On a un isom bifonctoriel en $F,x$
(si $F\in\deleted{\unsure{\in|\mathcal B}}\mathrm{Ob}\,\underline{\mathrm{Hom}}_{k\text{-}\mathrm{lin}}(C,\mathcal B)$)
\[
\begin{aligned}
 F(x)&\xrightarrow{\sim}\mathrm{Hom}_k^\bullet(\alpha^\bullet(x),F(C^\bullet))\\
 &\hspace{5.3em}\rotatebox{90}{$=$}\\
 &\hspace{1.3em}\mathrm{Hom}_k^\bullet(\alpha^\bullet(x),C_F^\bullet),
\end{aligned}
\]
[où on définit, pour un $K^\bullet\in\mathrm{Ob}\,\mathrm{Ab}_k^\bullet$
et $C^\bullet\in\mathrm{Ob}\,\mathcal B^\bullet_{\deleted{\unsure{k}}}$,
l'objet $\mathrm{Hom}_k^\bullet(K^\bullet,C^\bullet)$ de $\mathcal B$
(s'il existe) par l'isom fonctoriel en $b\in\mathrm{Ob}\,\mathcal B$
\[
 \mathrm{Hom}(\underset{\underset{\mathrm{Ob}}{\rotatebox{-90}{$\in$}}}{b},
 \mathrm{Hom}_k^\bullet(K^\bullet,C^\bullet))
 \simeq\mathrm{Hom}_{\mathcal B^\bullet}(b\otimes_k K^\bullet,C^\bullet),
\]
si les $K^i$ sont libres de type fini, $b\otimes_k K^\bullet$
est bien défini, donc $b\longmapsto\mathrm{Hom}_{\mathcal B^\bullet}(b\otimes_k K^\bullet,C^\bullet)$
l'est, et la question est si ce foncteur est représentable ; il se trouve qu'il l'est si
$K^\bullet$ est \deleted{\unsure{libre}} $=\in\mathrm{Ob}\,\deleted{\mathrm{Ab}_k}\mathcal K^\bullet$
[ Mais de façon générale, c'est le sous-foncteur de
% slanted
\margin{\unsure{remonter} avant énoncé prop 3}
\[
 \prod_i\underbrace{\mathrm{Hom}_k(K^i,C^i)}_{\in\mathrm{Ob}\,\mathcal B}
\]
noyau de
\[
 \prod_i\mathrm{Hom}_k(K^i,C^i)\mathrel{\underset{\deleted{\unsure{\longrightarrow}}}{\xrightarrow{\delta}}}\prod_i\mathrm{Hom}_k(K^i,C^{i+1})
\]
$(\delta(u^i)=du^i-u^{i+1}d)$ \unsure{...|on} Ça \deleted{est repr.} si
$\mathcal B$ \deleted{\unsure{abélienne}} \inserted{stable par Ker},
si les $\mathrm{Hom}_k(K^i,C^i)$ sont repr (p. ex. \deleted{\unsure{les}} $K^i$ de prés. finie) et les deux
% archive sheet number: 43
